% =====================================================================
% CSE 341 -- Concepts of Programming Languages
% Project D1 (Part 1) -- Design Specification for ANS
% Students: 220104004926 & 220104004931
% Compile with:  pdflatex 220104004926-220104004931_D1_P1.tex  (run twice)
% =====================================================================
\documentclass[11pt,a4paper]{article}

% ---------- Encoding, fonts, language ----------
\usepackage[utf8]{inputenc}
\usepackage[T1]{fontenc}
\usepackage{mathptmx}                    % Times serif (scalable Type 1)
\usepackage[scaled=0.92]{helvet}         % Helvetica sans
\usepackage{courier}                     % Courier monospace
\usepackage[english]{babel}
\usepackage{microtype}

% ---------- Page geometry ----------
\usepackage[a4paper,margin=2.0cm,top=2.0cm,bottom=2.2cm]{geometry}

% ---------- Colour, graphics, hyperlinks ----------
\usepackage[table]{xcolor}
\definecolor{ansblue}{HTML}{0B3D91}
\definecolor{ansrow}{HTML}{F4F6FB}
\definecolor{ansgray}{HTML}{F5F5F5}
\definecolor{ansborder}{HTML}{DDDDDD}
\definecolor{anscaption}{HTML}{555555}
\usepackage[colorlinks=true, linkcolor=ansblue, urlcolor=ansblue,
            citecolor=ansblue, pdfborder={0 0 0}]{hyperref}

% ---------- Typography ----------
\usepackage[explicit]{titlesec}
\titleformat{\section}
  {\color{ansblue}\Large\bfseries\sffamily}
  {\thesection}{0.6em}{#1}
\titleformat{\subsection}
  {\color{black}\normalsize\bfseries\sffamily}
  {\thesubsection}{0.6em}{#1}
\titlespacing*{\section}{0pt}{14pt}{6pt}
\titlespacing*{\subsection}{0pt}{10pt}{4pt}
\setlength{\parskip}{4pt plus 1pt}
\setlength{\parindent}{0pt}

% ---------- Tables ----------
\usepackage{array}
\usepackage{tabularx}
\usepackage{booktabs}
\usepackage{ragged2e}
\usepackage{makecell}
\usepackage{seqsplit}
\renewcommand{\arraystretch}{1.18}
\newcolumntype{Y}{>{\RaggedRight\arraybackslash}X}

% Header row helper (white text on ansblue)
\newcommand{\thd}[1]{\textbf{\color{white}#1}}

% ---------- Code listings ----------
\usepackage{listings}
\lstdefinelanguage{ANS}{
  morekeywords={let,function,return,for,in,if,elif,else,print,
                element,compound,reaction,reactants,products,
                formula,name,atomic_mass,valence,symbol,
                check,balanced,energy,exothermic,endothermic,
                molar_mass,compute_yield,mol},
  sensitive=true,
  morecomment=[l]{\#},
  morestring=[b]",
  morestring=[b]'
}
\lstdefinestyle{ans}{
  language=ANS,
  basicstyle=\ttfamily\footnotesize,
  keywordstyle=\bfseries\color{ansblue},
  commentstyle=\itshape\color{anscaption},
  stringstyle=\color{black!70!ansblue},
  numbers=none,
  frame=single,
  rulecolor=\color{ansborder},
  backgroundcolor=\color{ansgray},
  framesep=4pt,
  xleftmargin=4pt,
  xrightmargin=4pt,
  aboveskip=6pt,
  belowskip=6pt,
  showstringspaces=false,
  columns=fullflexible,
  keepspaces=true,
  breaklines=true,
  breakatwhitespace=true,
  upquote=true
}
\lstdefinestyle{ebnf}{
  language={},
  basicstyle=\ttfamily\footnotesize,
  keywordstyle={},
  commentstyle={},
  stringstyle={},
  numbers=none,
  frame=none,
  backgroundcolor=\color{white},
  aboveskip=4pt,
  belowskip=6pt,
  columns=fullflexible,
  keepspaces=true,
  breaklines=false,
  upquote=true,
  literate={->}{{$\rightarrow$}}2
}
\lstset{style=ans}

% Inline code shortcut
\newcommand{\code}[1]{\texttt{#1}}

% ---------- Header / footer ----------
\usepackage{fancyhdr}
\pagestyle{fancy}
\fancyhf{}
\renewcommand{\headrulewidth}{0pt}
\fancyfoot[L]{\footnotesize\color{anscaption}%
  CSE 341 --- ANS Design Specification (D1, Part 1) $\cdot$
  Students 220104004926 \& 220104004931}
\fancyfoot[R]{\footnotesize\color{anscaption}Page \thepage}

% ---------- PDF metadata ----------
\hypersetup{
  pdftitle  = {ANS -- Design Specification (D1, Part 1)},
  pdfauthor = {220104004926 and 220104004931},
  pdfsubject= {CSE 341 Project Part 1},
  pdfkeywords={DSL, chemistry, Sebesta, EBNF, lexical analysis}
}

% Suppress section numbers in TOC etc. -- not used
\setcounter{secnumdepth}{0}

% =====================================================================
\begin{document}

% ---------------- Title block ----------------
\begin{center}
  {\footnotesize\color{anscaption}
    CSE 341 --- Concepts of Programming Languages\par}
  \vspace{2pt}
  {\Huge\bfseries\sffamily\color{ansblue}
    ANS: A Domain-Specific Language for Chemistry\par}
  \vspace{4pt}
  {\normalsize Design Specification (D1) --- Part 1\par}
  {\footnotesize\color{anscaption}
    Students: 220104004926 \& 220104004931 \quad$\cdot$\quad
    Spring 2026 \quad$\cdot$\quad
    Gebze Technical University\par}
\end{center}
\vspace{6pt}

% =====================================================================
\section{4.1\quad Language Overview}

\textbf{Name.}\, ANS --- short for \textit{Atomic Notation Script}.

\textbf{Domain.}\, ANS is a small DSL for declaring chemical elements,
compounds, and reactions, verifying balance at compile time, and performing
stoichiometric queries such as molar mass and theoretical yield. It targets
undergraduate chemistry students and instructors who want to write
calculations the way they appear in textbooks rather than as bare numerical
code.

\subsection{Sample program}
\begin{lstlisting}
# Declare elements
element H { atomic_mass = 1.008,  valence = 1, symbol = "H" };
element O { atomic_mass = 15.999, valence = 2, symbol = "O" };
element C { atomic_mass = 12.011, valence = 4, symbol = "C" };

# Declare compounds
compound water   { formula = H2O, name = "Water"          };
compound methane { formula = CH4, name = "Methane"        };
compound co2     { formula = CO2, name = "Carbon dioxide" };
compound oxygen  { formula = O2,  name = "Oxygen"         };

# Variables and a heterogeneous list
let n        = 21.2;
let elements = [H, O, C];
let list     = [1, "two", 3.45, water, methane.formula, co2.name, H];

# A user-defined function
function molar_mass(compound) {
    let formula = compound.formula;
    let mass    = 0.0;
    for (let element in formula) {
        mass += element.atomic_mass * formula[element];
    }
    return mass;
};

# A reaction whose balance is checked at compile time
reaction combustion {
    reactants = CH4 + 2*O2,
    products  = CO2 + 2*H2O,
    check     = balanced,
    energy    = exothermic 890.3 kJ/mol
};

# Stoichiometric queries
molar_mass(water);
compute_yield(combustion, 5.0 mol, CH4);

# Conditional reporting
if compute_yield(combustion, 5.0 mol, CH4) > 0 {
    print "The yield of CH4 is positive, so combustion is possible.";
} elif compute_yield(combustion, 5.0 mol, CH4) == 0 {
    print "The yield is zero.";
} else {
    print "The yield is negative.";
}
\end{lstlisting}

\subsection{Sebesta evaluation-criteria justification (Ch.~1)}

Following Sebesta's Chapter~1 evaluation framework --- readability,
writability, reliability, and overall cost --- ANS deliberately prioritizes
\textbf{readability, writability, and reliability} while knowingly sacrificing
some \textbf{cost, generality, and orthogonality}. Readability is favoured
because the syntax mirrors chemistry vocabulary directly (\code{element},
\code{compound}, \code{reaction}, \code{reactants}, \code{products},
\code{balanced}) on top of a familiar C-like skeleton (\code{let}, braces,
semicolons, function calls, \code{if}/\code{elif}/\code{else}), so the meaning
of a program can be inferred from its surface form. Writability is favoured
because common domain tasks --- declaring an element, writing a balanced
equation as \code{CH4 + 2*O2 -> CO2 + 2*H2O}, or asking
\code{molar\_mass(water)} --- are first-class constructs rather than
user-coded library routines. Reliability is favoured through compile-time
validation: the lexer recognises formula and unit literals as their own
categories, the semantic checker verifies that every element symbol exists in
the periodic table and that every reaction is atomically balanced before any
value is computed, and units such as \code{mol} and \code{kJ/mol} are tracked
rather than discarded. The accepted trade-offs are the higher implementation
cost of formula parsing, unit handling, and balance checking, the loss of
generality (ANS is unsuited to non-chemistry tasks), and reduced orthogonality
(many fixed domain forms instead of a small set of freely combinable
primitives). Sebesta explicitly notes that language design is a process of
selecting trade-offs among these criteria; for a teaching-oriented chemistry
DSL the gains in domain-specific clarity and safety outweigh the loss of
general-purpose expressiveness.

% =====================================================================
\section{4.2\quad Lexical Structure}

ANS uses the standard five-category lexical split --- identifiers, literals,
keywords, operators, and separators --- that Sebesta describes as the output
of a finite-automaton-based lexical analyser (\S4.2). Two practical wrinkles
drive the design: (i)~chemical formulas like \code{H2O} must be
distinguishable from user-defined names, and (ii)~several chemistry words
(\code{reaction}, \code{compound}, \code{mol}) must still be usable as
ordinary parameter names inside function bodies. ANS solves (i) by case
(identifiers begin lowercase; element/formula tokens begin uppercase) and
(ii)~by treating the chemistry words as \textit{contextual} keywords that the
parser promotes only in the grammar positions where they are reserved.

\renewcommand{\tabularxcolumn}[1]{>{\RaggedRight\arraybackslash}m{#1}}
% A small helper so a regex column entry is set in monospace, scriptsize,
% and broken at any character if needed.
\newcommand{\rx}[1]{\scriptsize\ttfamily\seqsplit{#1}}
{\footnotesize
\rowcolors{2}{white}{ansrow}
\arrayrulecolor{ansborder}
\begin{tabularx}{\linewidth}{@{}>{\RaggedRight\arraybackslash}p{2.7cm}
                                >{\RaggedRight\arraybackslash}p{6.0cm}
                                Y@{}}
\arrayrulecolor{ansborder}\toprule
\rowcolor{ansblue}\thd{Category} & \thd{ANS meaning / examples} & \thd{Pattern (Python-style regex)}\\
\midrule
Identifiers
  & User-defined names: \code{water}, \code{methane}, \code{num}, \code{reactant\_mol}
  & \rx{[a-z\_][A-Za-z0-9\_]*}\\
Number literals
  & Integers, decimals, scientific: \code{20}, \code{1.008}, \code{6.02e23}
  & \rx{(?:0|[1-9][0-9]*)(?:\textbackslash.[0-9]+)?(?:[eE][+-]?[0-9]+)?}\\
String literals
  & \code{"Hydrogen"}, \code{"Water"}, \code{"Carbon dioxide"}
  & \rx{"(?:[\^{}"\textbackslash\textbackslash{}\textbackslash{}r\textbackslash{}n]|\textbackslash\textbackslash["\textbackslash\textbackslash{}nrt])*"}\\
Char literals
  & \code{'H'}, \code{'\textbackslash{}n'}
  & \rx{'(?:[\^{}'\textbackslash\textbackslash{}\textbackslash{}r\textbackslash{}n]|\textbackslash\textbackslash['\textbackslash\textbackslash{}nrt])'}\\
Formula literals
  & Element symbols / formulas: \code{H}, \code{O}, \code{H2O}, \code{CH4}, \code{CO2}, \code{O2}
  & \rx{(?:[A-Z][a-z]?(?:[1-9][0-9]*)?)+}\\
Unit literals
  & \code{mol}, \code{g}, \code{kg}, \code{g/mol}, \code{kg/mol}, \code{kJ}, \code{kJ/mol}
  & \rx{(?<![A-Za-z0-9\_])(?:kg/mol|g/mol|kJ/mol|mol|kg|g|kJ)(?![A-Za-z0-9\_])}\\
Keywords (hard)
  & \code{let}, \code{function}, \code{return}, \code{for}, \code{in}, \code{if}, \code{elif}, \code{else}, \code{print}
  & matched as identifier, then keyword-table lookup\\
Keywords (contextual)
  & \code{element}, \code{compound}, \code{reaction}, \code{reactants},
    \code{products}, \code{formula}, \code{name}, \code{atomic\_mass},
    \code{valence}, \code{symbol}, \code{check}, \code{balanced},
    \code{energy}, \code{exothermic}, \code{endothermic}
  & promoted to keywords only in their grammar positions\\
Operators
  & \code{=}\, \code{+}\, \code{-}\, \code{*}\, \code{/}\, \code{+=}\, \code{-=}\, \code{*=}\, \code{/=}\, \code{==}\, \code{!=}\, \code{<}\, \code{<=}\, \code{>}\, \code{>=}
  & \rx{==|!=|<=|>=|\textbackslash{}+=|-=|\textbackslash{}*=|/=|=|[+\textbackslash{}-*/<>]}\\
Separators
  & \code{\{}\, \code{\}}\, \code{(}\, \code{)}\, \code{[}\, \code{]}\, \code{,}\, \code{;}\, \code{.}
  & \rx{[\{\}()\textbackslash{}[\textbackslash{}],;.]}\\
Comments
  & \code{\#} line comment to end of line
  & \rx{\#[\^{}\textbackslash{}r\textbackslash{}n]*}\\
Whitespace
  & spaces, tabs, newlines (separator only, not a token)
  & \rx{[\textbackslash{} \textbackslash{}t\textbackslash{}f\textbackslash{}r\textbackslash{}n]+}\\
\bottomrule
\end{tabularx}}
\par\vspace{4pt}

\textbf{Lexer ordering rules.}\, Match the longer pattern first within each
group: \code{kJ/mol} before \code{kJ}; \code{g/mol} before \code{g}; compound
assignments such as \code{+=} before \code{+}; and \code{==}, \code{<=},
\code{>=}, \code{!=} before their single-character forms. Identifiers are
matched first, then promoted to \textit{KEYWORD} if their lexeme is in the
reserved-word table --- Sebesta's standard technique for names vs. reserved
words (\S5.1).

% =====================================================================
\newpage
\section{4.3\quad Syntax (EBNF)}

The grammar below is written in Sebesta-style EBNF (\S3.1--3.2): non-terminals
appear in angle brackets, terminals in double quotes, \code{\{ X \}} means
zero or more repetitions of X, \code{[ X ]} means optional, and \code{|}
separates alternatives. The grammar is unambiguous for the implemented
subset; ambiguity-resolution decisions are listed at the end of this section.

\begin{lstlisting}[style=ebnf]
(* ============================================================
   ANS Grammar (Part 1) -- Sebesta-style EBNF
   ============================================================ *)

(* ---------- Program structure ---------- *)
<program>           ::= { <top_level_statement> } EOF

<top_level_statement> ::= <variable_declaration> ";"
                        | <assignment_statement>  ";"
                        | <element_declaration>   ";"
                        | <compound_declaration>  ";"
                        | <reaction_declaration>  ";"
                        | <function_declaration>
                        | <if_statement>
                        | <expression_statement>  ";"

(* ---------- Blocks and statements ---------- *)
<block>             ::= "{" { <statement> } "}"

<statement>         ::= <variable_declaration>           ";"
                      | <assignment_statement>           ";"
                      | <augmented_assignment_statement> ";"
                      | <return_statement>               ";"
                      | <for_statement>
                      | <if_statement>
                      | <print_statement>                ";"
                      | <expression_statement>           ";"

<variable_declaration>  ::= "let" <identifier> "=" <expression>
<assignment_statement>  ::= <assignable> "=" <expression>
<augmented_assignment_statement>
                        ::= <assignable> <aug_op> <expression>
<aug_op>                ::= "+=" | "-=" | "*=" | "/="
<return_statement>      ::= "return" [ <expression> ]
<print_statement>       ::= "print" <expression>
<expression_statement>  ::= <expression>

(* ---------- Functions and loops ---------- *)
<function_declaration> ::= "function" <identifier>
                           "(" [ <parameter_list> ] ")" <block> ";"
<parameter_list>       ::= <identifier> { "," <identifier> }

<for_statement>        ::= "for" "(" "let" <identifier> "in" <expression> ")"
                           <block>

(* ---------- Conditionals (no dangling-else: braces required) ---------- *)
<if_statement>         ::= "if" <expression> <block>
                           { "elif" <expression> <block> }
                           [ "else" <block> ]

(* ---------- Domain-specific declarations ---------- *)
<element_declaration>  ::= "element" <element_symbol>
                           "{" <element_field_list> "}"
<element_symbol>       ::= <formula_literal>
<element_field_list>   ::= <element_field> { "," <element_field> } [ "," ]
<element_field>        ::= "atomic_mass" "=" <number_literal>
                         | "valence"     "=" <number_literal>
                         | "symbol"      "=" <string_literal>

<compound_declaration> ::= "compound" <identifier>
                           "{" <compound_field_list> "}"
<compound_field_list>  ::= <compound_field> { "," <compound_field> } [ "," ]
<compound_field>       ::= "formula" "=" <formula_literal>
                         | "name"    "=" <string_literal>

<reaction_declaration> ::= "reaction" <identifier>
                           "{" <reaction_field_list> "}"
<reaction_field_list>  ::= <reaction_field> { "," <reaction_field> } [ "," ]
<reaction_field>       ::= "reactants" "=" <reaction_expression>
                         | "products"  "=" <reaction_expression>
                         | "check"     "=" <balance_check>
                         | "energy"    "=" <energy_expression>
<balance_check>        ::= "balanced"
<energy_expression>    ::= <energy_kind> <quantity_literal>
<energy_kind>          ::= "exothermic" | "endothermic"

(* ---------- Reaction expressions (own grammar, separate from arithmetic) ---------- *)
<reaction_expression>  ::= <reaction_term> { "+" <reaction_term> }
<reaction_term>        ::= <reaction_species>
                         | <coefficient> "*" <reaction_species>
<coefficient>          ::= <number_literal>
<reaction_species>     ::= <formula_literal> | <identifier>

(* ---------- General expressions (precedence encoded in the grammar) ---------- *)
<expression>                ::= <or_expression>
<or_expression>             ::= <and_expression>
                                { "||" <and_expression> }
<and_expression>            ::= <equality_expression>
                                { "&&" <equality_expression> }
<equality_expression>       ::= <relational_expression>
                                { ("==" | "!=") <relational_expression> }
<relational_expression>     ::= <additive_expression>
                                { ("<" | "<=" | ">" | ">=") <additive_expression> }
<additive_expression>       ::= <multiplicative_expression>
                                { ("+" | "-") <multiplicative_expression> }
<multiplicative_expression> ::= <unary_expression>
                                { ("*" | "/") <unary_expression> }
<unary_expression>          ::= [ "-" ] <postfix_expression>
<postfix_expression>        ::= <primary_expression> { <postfix_op> }
<postfix_op>                ::= "." <identifier>
                              | "[" <expression> "]"
                              | "(" [ <argument_list> ] ")"

<primary_expression>        ::= <quantity_literal>
                              | <number_literal>
                              | <string_literal>
                              | <char_literal>
                              | <formula_literal>
                              | <list_literal>
                              | <identifier>
                              | "(" <expression> ")"

<argument_list>             ::= <expression> { "," <expression> }
<quantity_literal>          ::= <number_literal> <unit_literal>
<list_literal>              ::= "[" [ <expression> { "," <expression> } ] "]"

(* ---------- Assignable LHS (subset of postfix; no calls allowed) ---------- *)
<assignable>          ::= <identifier> { <assignable_postfix> }
<assignable_postfix>  ::= "." <identifier>
                        | "[" <expression> "]"

(* ---------- Lexical (terminals from the lexer) ---------- *)
<identifier>      ::= IDENTIFIER       (* [a-z_][A-Za-z0-9_]*           *)
<number_literal>  ::= NUMBER_LITERAL
<string_literal>  ::= STRING_LITERAL
<char_literal>    ::= CHAR_LITERAL
<formula_literal> ::= FORMULA_LITERAL  (* (?:[A-Z][a-z]?(?:[1-9][0-9]*)?)+ *)
<unit_literal>    ::= UNIT_LITERAL
\end{lstlisting}

\subsection{Ambiguity-resolution annotations}

{\footnotesize
\rowcolors{2}{white}{ansrow}
\arrayrulecolor{ansborder}
\begin{tabularx}{\linewidth}{@{}>{\RaggedRight\arraybackslash}p{4.0cm} Y@{}}
\arrayrulecolor{ansborder}\toprule
\rowcolor{ansblue}\thd{Source of ambiguity} & \thd{How ANS resolves it}\\
\midrule
Operator precedence
  & Encoded by the cascade
    \textit{or} $\to$ \textit{and} $\to$ equality $\to$ relational $\to$
    additive $\to$ multiplicative $\to$ unary $\to$ postfix $\to$ primary.
    Therefore \code{a + b * c} parses as \code{a + (b * c)} and
    \code{a < b + c} parses as \code{a < (b + c)}. All binary operators are
    left-associative because each level repeats with \code{\{ ... \}}.\\
Reaction vs. arithmetic +/*
  & Reactions have their own non-terminal \code{<reaction\_expression>},
    separate from \code{<expression>}. Inside a reaction, \code{+} separates
    chemical species and \code{*} binds a coefficient to a formula, so
    \code{CH4 + 2*O2} parses unambiguously as \code{CH4 + (2*O2)}.\\
Stoichiometric coefficient vs. subscript
  & ANS requires the explicit form \code{2*O2} rather than \code{2O2}. This
    removes the ambiguity between a coefficient and a leading digit of a
    (non-existent) formula token, and keeps \code{FORMULA\_LITERAL} purely
    uppercase-led.\\
Formula literal vs. identifier
  & Identifiers begin with a lowercase letter or underscore; formula literals
    begin with an uppercase letter. Thus \code{water} is always an identifier
    and \code{H2O} is always a formula. A bare \code{H} is lexically a
    \code{FORMULA\_LITERAL}; the parser interprets it as an
    \code{<element\_symbol>} only when it follows the keyword \code{element}.\\
Contextual keywords
  & \code{reaction}, \code{compound}, \code{mol}, \code{formula}, etc.\ are
    recognised as keywords only in the grammar positions where they appear
    above. Outside those positions they remain identifiers, which is why
    \code{function compute\_yield(reaction, mol, compound) \{ ... \}} is
    legal.\\
Dangling-else
  & ANS requires braces around every then/elif/else body
    (see \code{<if\_statement>}). The classical dangling-else ambiguity
    therefore cannot arise: every \code{else} attaches to a uniquely
    delimited \code{<block>}.\\
Function call vs. assignment target
  & \code{<assignable>} deliberately omits the call postfix
    \code{'(' ... ')'}. Therefore \code{f(x) = e} is a syntax error,
    removing the ambiguity between calling a function and assigning to its
    result.\\
Field order in declarations
  & \code{<element\_field\_list>}, \code{<compound\_field\_list>}, and
    \code{<reaction\_field\_list>} accept fields in any order syntactically;
    the semantic checker enforces that each required field appears exactly
    once. This keeps the grammar LL(1) while still rejecting duplicate or
    missing fields with a precise error message.\\
\bottomrule
\end{tabularx}}

% =====================================================================
\newpage
\section{4.6\quad Names, Binding, Scope, Lifetime}

\subsection{Rules for legal identifiers (Sebesta \S5.1)}

An ANS identifier matches \code{[a-z\_][A-Za-z0-9\_]*}: it begins with a
lowercase letter or underscore and continues with letters, digits, or
underscores. ANS is \textbf{case-sensitive}, so \code{water} and \code{Water}
are distinct names. Tokens that begin with an uppercase letter (e.g.\
\code{H}, \code{O2}, \code{H2O}) are \textit{formula literals}, never
identifiers, which keeps chemistry notation visually separated from program
names. The hard reserved words \code{let, function, return, for, in, if,
elif, else, print} may never be used as identifiers. The chemistry words
\code{element, compound, reaction, reactants, products, formula, name,
atomic\_mass, valence, symbol, check, balanced, energy, exothermic,
endothermic} are \textit{contextual} keywords in the sense of Sebesta's
discussion of keywords vs.\ reserved words: they are reserved only where the
grammar expects them, so
\code{function compute\_yield(reaction, mol, compound) \{ ... \}} is a legal
declaration. A name may not be declared twice in the same scope; a name in
an inner scope may shadow an outer name, and the type checker emits a warning
when this shadows a top-level chemistry declaration.

\subsection{Bindings: compile time vs.\ run time (Sebesta \S5.2)}

{\footnotesize
\rowcolors{2}{white}{ansrow}
\arrayrulecolor{ansborder}
\begin{tabularx}{\linewidth}{@{}>{\RaggedRight\arraybackslash}p{5.4cm}
                               >{\RaggedRight\arraybackslash}p{4.0cm}
                               Y@{}}
\arrayrulecolor{ansborder}\toprule
\rowcolor{ansblue}\thd{Binding} & \thd{When} & \thd{Example}\\
\midrule
Operator symbol $\to$ meaning
  & Language-design time
  & \code{+}, \code{*}, \code{==}, \code{\&\&}\\
Lexeme $\to$ token category
  & Lex / compile time
  & \code{water} $\to$ IDENTIFIER, \code{H2O} $\to$ FORMULA\_LITERAL\\
Identifier $\to$ declaration
  & Compile time (static scope walk)
  & \code{mass} refers to its enclosing \code{let}\\
Variable $\to$ inferred type
  & Compile time
  & \code{let n = 21.2;} $\Rightarrow$ \code{n : Float}\\
Element / compound / reaction name $\to$ declaration
  & Compile time
  & \code{water}, \code{combustion}\\
Formula validity, balance check, unit compatibility
  & Compile time
  & \code{check = balanced} verified statically\\
Parameter $\to$ argument value
  & Run time
  & \code{compute\_yield(combustion, 5.0 mol, CH4)}\\
Local variable $\to$ memory cell
  & Run time (on entry)
  & stack-allocated when block is entered\\
Variable $\to$ value
  & Run time
  & \code{mass += element.atomic\_mass * formula[element]}\\
Function name $\to$ return value
  & Run time
  & \code{return product\_mol;}\\
\bottomrule
\end{tabularx}}
\par\vspace{2pt}

ANS therefore uses mostly \textbf{static} bindings for names, types, and
chemistry validation, and \textbf{dynamic} bindings for storage and values
--- exactly the mix Sebesta recommends when reliability is more important
than late flexibility (\S5.4.3).

\subsection{Scoping rule: static (lexical) scoping (Sebesta \S5.5)}

ANS uses \textbf{static, lexical scoping}. The reference of every name is
determined from the program text alone: a name resolves to the nearest
enclosing declaration in the source --- block, loop body, function body,
then enclosing function (if nested), then top level. The reasoning follows
Sebesta's argument that static scope makes both the human reader and the
compiler able to determine the meaning of every name without simulating
execution.

\textbf{What would break under dynamic scoping.}\, Consider:

\begin{lstlisting}
let water = methane;          # 'water' shadowed at top level

function f() {
    return molar_mass(water);  # what does 'water' refer to?
}

function g() {
    let water = co2;
    return f();                # caller introduces a local 'water'
}
\end{lstlisting}

Under static scoping \code{f} always uses the top-level \code{water} visible
where it was \textit{defined}, so \code{g()} still computes the molar mass of
methane. Under dynamic scoping \code{f} would instead capture \code{g}'s
local \code{water} and silently switch to CO\textsubscript{2}. For a chemistry
DSL this is intolerable: a stoichiometric calculation must not depend on the
call chain. Dynamic scoping would also defeat the compile-time
formula-validity and balance checks, because the type checker could no longer
determine which compound a function will read at run time. ANS therefore
commits to static scoping; the small loss of run-time flexibility is exactly
the trade-off named in \S4.1.

\subsection{Lifetimes (Sebesta \S5.4.3)}

{\footnotesize
\rowcolors{2}{white}{ansrow}
\arrayrulecolor{ansborder}
\begin{tabularx}{\linewidth}{@{}>{\RaggedRight\arraybackslash}p{5.0cm}
                               >{\RaggedRight\arraybackslash}p{4.0cm}
                               Y@{}}
\arrayrulecolor{ansborder}\toprule
\rowcolor{ansblue}\thd{Entity} & \thd{Lifetime category} & \thd{Notes}\\
\midrule
Top-level element / compound / reaction declarations
  & Static (program lifetime)
  & Allocated once at program start; visible to every later statement and
    function.\\
Top-level \code{let} variables
  & Static (program lifetime)
  & Single binding from the \code{let} until program termination.\\
Function parameters
  & Stack-dynamic
  & Allocated on call; deallocated on return. Supports recursion.\\
\code{let} variables inside a function or block
  & Stack-dynamic
  & Allocated when the declaration is elaborated; freed when the enclosing
    block exits.\\
Loop variable in \code{for (let x in ...)}
  & Stack-dynamic, scoped to the loop
  & A fresh binding per iteration, so closures over \code{x} do not alias.\\
List values, compound objects, reaction objects
  & Heap-managed by the implementation
  & Not user-visible: ANS has no \code{new} / \code{delete} and no pointers,
    so heap management is automatic, not ``explicit heap-dynamic'' in
    Sebesta's sense.\\
\bottomrule
\end{tabularx}}
\par\vspace{2pt}

Local variables are \textbf{stack-dynamic} rather than static so that
repeated calls to \code{molar\_mass} or \code{compute\_yield} cannot leak
state from one invocation to the next, and so that recursion (anticipated for
Part~2's reaction search and balance solver) is supported. ANS deliberately
exposes no \textbf{explicit heap-dynamic} construct: there are no pointers,
no \code{new}/\code{delete}, and no aliasing. This costs some flexibility but
eliminates an entire class of unreliability errors --- a trade-off that is
consistent with the readability/reliability prioritisation declared in
\S4.1.

\vspace{6pt}
{\footnotesize\itshape\color{anscaption}
Part~2 will extend this document with \S4.4 Semantics, \S4.5 Type System,
\S4.7 Expressions \& Assignment, and \S4.8 Design Rationale, all consistent
with the decisions stated above.\par}

\end{document}
